\scanpage{55}
% E: text|left upper, corollary|Dans tous les cas
\uline{Corollaire} Dans tous les cas, la distribution des signes \(\varepsilon_{\unsure{a|\cdot}}(F)\) détermine celle des \(\varepsilon_{\unsure{a'|\cdot}}(F')\). \deleted{C'est clair \unsure{et}}

% E: text|left upper-middle, proof|C'est clair d'abord
C'est clair d'abord dans les cas \(1^\circ\) et \(3^\circ\), où on sait laquelle des faces \(F,F'\) est la spécialisée de l'autre ; \deleted{et par suite}\inserted{Mais} dans le cas\inserted{``symétrique''} \((2^\circ)\), \deleted{\unsure{comme on connaît les signes \(\varepsilon(F)\), on}} On voit que \deleted{les} deux\inserted{grands} côtés extérieurs correspondants pour \(F,F'\) ont $=$ \uline{signe} (c'est tjrs le cas, y compris dans les cas \(1^\circ\) et \(2^\circ\)), et pour les côtés qui \unsure{restent}\inserted{\unsure{en}} par dessus, elles sont (pour \(F,F'\) \unsure{par \(\sigma\)}) toutes de $=$ signe (\uline{NB} leur nb, suivant les cas, est \((1,0),(1,1)\) ou \((1,2)\)) et les signes en question pour \(F\) et pour \(F'\) sont opposés.

% E: text|left bottom, second corollary|Quand on connaît les signes
\uline{Corollaire} Quand on connaît les signes \(\varepsilon_F(\unsure{a})\) pour \uline{une} face \(F\), on les connaît pour \uline{tous} (grâce à la connexité de \(X\))

% E: text|right upper, qualification|J'ai supposé implicitement
J'ai supposé implicitement \deleted{que} dans la vérification de la proposition que ni \(F\), ni \(F'\) n'est une face exceptionnelle (i.e. ne correspondant à une position \unsure{ultime|critique} exceptionnelle, i.e. ``supercritique''). Il faudrait que j'examine séparément le cas de telles faces.

% E: text|right middle, exceptional face|Pour une telle face
Pour une telle face, il n'y a \unsure{que|pas} quatre grands côtés, et la distribution des signes est connue --- Cela semble donner une façon simple de déduire les signes\inserted{ou les rôles} de la seule structure de carte de \((X,K,S_K)\), sans avoir à utiliser les arcs critiques, et encore moins \deleted{les} le \deleted{\unsure{complexe}} graphe \(L\subset X\).
\mspage{54}
